\subsection{Hadamard Transform on Multiple Qubits}

Applying the Hadamard gate to every qubit is a useful tool in quantum computing, as it creates a superposition of all possible states. Now, let's see how this is done.

\highspace
We already know that applying the Hadamard gate to a single qubit in the state $\ket{0}$ gives us:
\begin{equation*}
    H\ket{0} = \frac{1}{\sqrt{2}}\left(\ket{0} + \ket{1}\right)
\end{equation*}
This means that $H$ creates a \textbf{balanced superposition} of the two basis states $\ket{0}$ and $\ket{1}$.

\highspace
Now, if we apply this idea to two qubits, we can apply the Hadamard gate to each qubit individually. This is done using the tensor product of the Hadamard gates, for example:
\begin{equation*}
    \left(H \otimes H\right) \ket{00} = H\ket{0} \otimes H\ket{0}
\end{equation*}
So it is equivalent to applying the Hadamard gate to each qubit separately. Expanding this, we get:
\begin{align*}
    \left(H \otimes H\right) \ket{00} &= H\ket{0} \otimes H\ket{0} \\
    &= \left(\frac{1}{\sqrt{2}}\left(\ket{0} + \ket{1}\right)\right) \otimes \left(\frac{1}{\sqrt{2}}\left(\ket{0} + \ket{1}\right)\right) \\
    &= \frac{1}{2} \left(\ket{0} + \ket{1}\right) \otimes \left(\ket{0} + \ket{1}\right) \\
    &= \frac{1}{2} \left(\ket{00} + \ket{01} + \ket{10} + \ket{11}\right)
\end{align*}
This means that \hl{applying Hadamard gate in parallel} \textbf{generates all possible combinations of the basis states} for the two qubits, creating a superposition of all four possible states.

\highspace
In general, for $n$ qubits, \hl{applying the Hadamard gate to each qubit will create a superposition of all $2^n$ possible states.} The resulting state can be expressed as:
\begin{equation}\label{eq:hadamard-transform-multiple-qubits}
    \left(H^{\otimes n}\right) \ket{0}^{\otimes n} = \frac{1}{2^{\frac{n}{2}}} \sum_{k = 0}^{2^n - 1} \ket{k} = \frac{1}{\sqrt{2^n}} \sum_{k = 0}^{2^n - 1} \ket{k}
\end{equation}
Where $\ket{k}$ represents the basis states of the $n$-qubit system, and the \textbf{sum runs over all possible combinations of the basis states}.

\begin{figure}[!htp]
    \centering
    \begin{quantikz}
        \lstick{$\ket{0}$} & \gate{H} & \rstick[2]{$\left(H \otimes H\right)\ket{00} = \dfrac{1}{2}(\ket{00} + \ket{01} + \ket{10} + \ket{11})$} \\
        \lstick{$\ket{0}$} & \gate{H} & 
    \end{quantikz}
\end{figure}

\noindent
In summary, \textbf{Hadamard on all qubits spreads the state over the entire space of possible states}, creating a uniform superposition that is fundamental in many quantum algorithms.